\begin{frame}{실습: 세 문제를 스스로 정식화해보기 (정답)}
  \small
  각 시나리오에 대해 $x$(입력 특징), $y$(출력), $h$ 의 대략적 형태(직선?
  규칙 기반? ``이번 학기에 배울 것'')를 노트에 적어보기.
  \begin{tabular}{@{}p{0.34\textwidth}p{0.6\textwidth}@{}}
    \toprule
    시나리오 & 힌트 (x / y / 갈래) \\
    \midrule
    \textbf{1. 신용카드 이상거래 탐지}\\ 매 결제가 사기인지를 즉시 판단 &
    $x$=결제 금액·시간·가맹점 종류, \; $y$=사기/정상(범주) $\to$ \textbf{분류} \\
    \midrule
    \textbf{2. 뉴스 기사 주제 자동 분류}\\ 정치/경제/스포츠/문화 태그 &
    $x$=기사 텍스트(나중에 단어 빈도·임베딩으로), \; $y$=4개 범주 중
    하나 $\to$ \textbf{분류(다중 클래스)} \\
    \midrule
    \textbf{3. 공장 설비 잔여수명 예측}\\ 진동·온도·가동시간 센서 &
    $x$=센서 측정값들, \; $y$=남은 일수(숫자) $\to$ \textbf{회귀} \\
    \bottomrule
  \end{tabular}
  \vskip0.6em
  세 문제 모두 ``\textbf{정답 라벨이 있고 $y$ 가 범주냐 숫자냐}''라는
  같은 질문 하나로 갈린다 --- 이번 학기 내내 새 문제를 만날 때마다
  가장 먼저 꺼내 쓸 도구.
\end{frame}
